% ==============================================================================
% SEZIONE 5: LA TRANSIZIONE DIGITALE: GRIDS E SMART METERING
% ==============================================================================

\section{Transizione Digitale}

\begin{frame}{Infrastrutture di Rete: La Centralità delle Grids}
    La digitalizzazione delle reti rappresenta l'architrave su cui Enel poggia la propria efficienza operativa e la nuova strategia di relazione con il cliente.
    \vspace{0.2cm}
    
    \note{La digitalizzazione delle reti trasforma le grids da asset fisici a vettori di efficienza operativa.}

    \begin{block}{Pianificazione degli Investimenti (Capex)}
        Il Gruppo ha accelerato massicciamente nel segmento \textbf{Grids}:
        \begin{itemize}
            \item \textbf{Investimento Totale:} €47,5 miliardi destinati alla resilienza.
            \item \textbf{Obiettivo:} Accogliere fonti rinnovabili intermittenti e stabilizzare il carico.
        \end{itemize}
    \end{block}
    \vspace{0.1cm}
    \pause
    \note{Un piano capex di quarantasette miliardi e mezzo punta ad aumentare la resilienza per assorbire la generazione rinnovabile intermittente.}

    \begin{block}{Sostentamento dei Margini}
        Con l'aumento della \textbf{RAB} in Italia (+10\%) e Iberia, Enel trasforma il capitale investito in flussi di cassa regolati, sicuri e prevedibili.
    \end{block}
    \note{La crescita della RAB, in aumento del dieci per cento in Italia, converte gli investimenti in flussi di cassa regolati e prevedibili.}
\end{frame}


\begin{frame}{Smart Meter 2G e l'Ecosistema "Sticky"}
    Il contatore di seconda generazione (2G) è il dispositivo abilitante della strategia di \textit{customer empowerment} e della stabilità del credito.
    \vspace{0.2cm}

    \begin{columns}[T]
        \begin{column}{0.48\textwidth}
            \begin{block}{Specifiche Tecniche 2G}
                Monitoraggio in tempo reale e campionamento dei dati ogni \textbf{15 minuti}. Abilita il passaggio da consumatore a \textbf{prosumer}.
            \end{block}
        \end{column}
        \pause
        \note{I contatori di seconda generazione offrono dati ogni quindici minuti, abilitando l'evoluzione dell'utente in prosumer attivo.}
        
        \begin{column}{0.48\textwidth}
            \begin{block}{E-BOX: Digital Saving Bank}
                Valorizzazione dei kWh non consumati. Crea un ecosistema premiante che riduce drasticamente il \textbf{Bad Debt}.
            \end{block}
        \end{column}
    \end{columns}
    \vspace{0.3cm}
    \pause
    \note{Il sistema E-BOX è il proof of concept di un ecosistema sticky che fidelizza il cliente valorizzando i risparmi.}
    
    \begin{alertblock}{Impatto sulla Qualità del Credito}
        La tecnologia crea un sistema "premiante" che incentiva la regolarità dei pagamenti e riduce la perdita di clienti.
    \end{alertblock}
    \note{L'automazione digitale riduce il bad debt incentivando i pagamenti regolari, migliorando la qualità del credito.}
\end{frame}
